\Askhsh{20943}{4}
\begin{ylh}{1}
\item 3.2 Βασικές Τριγωνομετρικές Ταυτότητες
\item 3.3 Αναγωγή στο 1o Τεταρτημόριο
\item 3.4 Οι τριγωνομετρικές συναρτήσεις
\item 4.3 Πολυωνυμικές εξισώσεις και ανισώσεις
\end{ylh}
Δίνεται γωνία $x$ με $\frac{3\pi}{2} < x < 2\pi$ και οι παραστάσεις:
$Α = \hm^{2}(\pi - x) + \hm^{2}(\pi + x) + \syn^{2}( - x)$,
$B = \frac{\hm  x}{1 + \syn  x} + \frac{1 + \syn  x}{\hm  x}$ .
\begin{alist}
\item Να αποδείξετε ότι$\ \ \ Α = \hm^{2}x + 1.$
\monades{08}
\item Να απλοποιήσετε την παράσταση $B$.
\monades{08}
\item Να εξετάσετε αν υπάρχει γωνία x για την οποία οι παραστάσεις Α και Β να είναι ίσες.
\monades{09}
\end{alist}

